L'objectif du module n'est pas d'accumuler des formules, mais d'acquérir une méthode : partir d'une situation réelle, identifier les bonnes variables, expliciter les hypothèses, effectuer un calcul avec la loi de probabilité adéquate, puis interpréter le résultat.

\section*{Objectifs du module}

À l'issue du cours, vous devez être capable de :
\begin{itemize}
    \item traduire une situation concrète en langage probabiliste 
    \item choisir un outil adapté parmi les lois et méthodes usuelles 
    \item calculer une probabilité, une espérance, une variance ou un intervalle de confiance 
    \item commenter le résultat en français simple dans un contexte donné 
    \item discuter la pertinence d'un modèle et les limites d'une conclusion.
\end{itemize}

\section*{Esprit du poly}

Ce document est un \textbf{support d'usage}. Il contient :
\begin{itemize}
    \item les définitions à connaître 
    \item les formules utiles 
    \item des méthodes de résolution 
    \item des exemples appliqués au domaine de la construction et du contrôle
\end{itemize}

\section*{Fils rouges}

Tout au long du poly, plusieurs situations reviennent régulièrement :
\begin{itemize}
    \item le \textbf{contrôle qualité du béton} et des éprouvettes ;
    \item la \textbf{planification de chantier} sous aléa météorologique ;
    \item le \textbf{contrôle de production} sur des lots de pièces ;
    \item l'\textbf{estimation statistique} à partir d'un échantillon de mesures.
\end{itemize}
